\scanpage{40}
\mspage{39}
\mspage{-4}
\uline{Professeur Nicolaas H. Kuiper}: \textit{Variétés Hilbertiennes}

\uline{Programme}

Variétés Hilbertiennes; aspects géométriques.

Les sujets qui seront traités sont les suivants:

\uline{Introduction}: Définition d'une variété $X$, modelée sur un espace de Banach, de classe de différentiabilité $C^r$. (Le livre de Serge Lang: ``Introduction to differentiable manifolds'' sera supposé connu de l'auditoire). (There is a cheaper french version of the book).

Théorème de Bessaga: L'espace de Hilbert $H$ est (analytiquement) difféomorphe à $H-\{0\}$.

Construction d'un plongement fermé de $X$ dans $B$. Par exemple pour $B=c_o\oplus\unsure{l|1}_p$. Existence d'une norme $C^\infty$ sur l'espace , $c_o$ des suites convergeant vers $0$.

Le groupe linéaire d'un espace de Hilbert est contractile.

Le théorème d'isotopie ambiante des voisinages tubulaires.

La version, en dimension infinie, du théorème d'isotopie ambiante de Hirsch et quelques lemmes simples de transversalité.

Fonctions non dégénérées sur une variété hilbertienne, condition $C$ de Palais-Smale. Forme canonique d'une telle fonction au voisinage d'un point critique. Théorème d'existence de Nicole Moulis. Décomposition en anses. Théorie de Morse. Stabilité d'après Palais: la variété Hilbertienne $X$ est difféomorphe à $X\,x\,H$. Enoncé du théorème de plongement ouvert de Eells et Elworthy.

\margin{$\ldots/$}
